\section{Definition of Observer}
\label{chap:definition-of-observer}

\subsection*{Axioms}

The argument begins with four axioms that serve as the premises for the entire derivation.

\textbf{Axiom 1: Genetic Encoding of Subjective Experience}
DNA is the blueprint of a conscious, pain-sensitive human.

\[
A_1: \text{Under the observed physical laws }(P),\ \text{DNA encodes the information necessary to construct systems capable of subjective experience }(S).
\]

(Where $D$ is the DNA/Genome, $H$ is the human organism, and $S$ is subjective experience.)

\textbf{Axiom 2: Physicality and Axiomatic Law}
DNA, the human organism, and the surrounding environment are composed solely of ordinary physical matter governed by physical laws ($P$).
\[
A_2: (H, U) \in \text{Axiomatic System}(P, D)
\]
(Where $U$ denotes the physical environment required to support $H$, and $P$ denotes the governing physical laws.)

\textbf{Axiom 3: Generalized Church-Turing Thesis Holds}

Any physically realizable process that can be described mathematically can, in principle, be simulated by a computational system.


\textbf{Axiom 4: Pain Affects Behavior}
A human with subjective experience ($S$, e.g., pain) behaves differently than the same human without it.
\[
H + S \neq H
\]
If a simulation $H'$ is identical to $H$ in behavior, it must also include $S$.

\subsection*{Deduction 1: Substrate Independence and Time as an Internal Property of the Observer}

\textbf{The Optimization Argument}

Consider a DNA simulation running on a computer, denoted by $T_{\text{alg}}$, consisting of code (the laws of physics) and data (the state of the simulated universe).
One can gradually optimize the code by introducing lookup tables, eventually replacing all computation with a static, precomputed dataset, denoted by $T_{\text{data}}$.

\begin{itemize}
    \item Let $E_{\text{int}}$ be Alice's experience of time and pain (the internal state transitions).
    \item Let $E_{\text{ext}}$ be the computer's external runtime (number of CPU cycles).
\end{itemize}

\textbf{Premise:} Code optimization changes $E_{\text{ext}}$ but preserves $E_{\text{int}}$.
\[
\text{Optimization}(T_{\text{alg}}) \rightarrow T_{\text{data}} \implies E_{\text{int}}(T_{\text{alg}}) \equiv E_{\text{int}}(T_{\text{data}})
\]

In this limit, the external runtime \(E_{\text{ext}}\) becomes zero in the sense that no state transitions are executed; the complete execution trace exists as static data.
One can therefore ask: does Alice's consciousness still persist in $T_{\text{data}}$?

If consciousness were to cease in $T_{\text{data}}$, it would necessitates a minimum code/data ratio for subjective experience.
This minimum ratio would be a new, non-physical constant imposed on $A_2$, leading to a contradiction.

\begin{quote}
  \textbf{Conclusion:} Consciousness can emerge from pure static data.
  Time and subjective experience ($S$) must emerge solely from information.
\end{quote}

\textbf{The Multi-threaded Argument}

Consider a multi-threaded computer running two DNA simulations, $\tau_A$ (Alice) and $\tau_B$ (Bob), concurrently, with a minimal time slice of one CPU cycle per thread.
The execution trace of the computer is then an interleaved sequence of segments drawn from both $\tau_A$ and $\tau_B$.
As the number of threads (simulated observers) increases without bound, the resulting execution trace asymptotically approaches white noise.
Consequently, each bit originating from $\tau_A$ are separated by increasingly large intervals of unrelated data.

Is Alice still conscious in this limit? If one maintains that $E_{\text{int}}$ vanishes as the number of threads increases, one must define a
specific thread density or bit-contiguity threshold required for subjectivity. Such a threshold would constitute a new physical constant,
which contradicts the completeness of $A_2$.

\begin{quote}
\textbf{Conclusion:} Conscious experience can arise from static white noise.
\end{quote}

\subsection*{Quantum Computers}
The preceding arguments considered simulations on classical computers. The same conclusions apply to quantum computers. Quantum computation provides an alternative physical realization of the same formal state-transition structure. Since the axioms establish substrate independence, any physically realizable computation — whether classical or quantum — can support subjective experience ($S$) provided it implements the relevant informational dynamics. Even if quantum computers are used, the optimization and multi-threading limits still apply, leading to the same static, timeless informational ontology.

\textbf{Numerical Precision}
Physical simulations are necessarily approximate. Let $\epsilon$ denote the numerical precision (or discretization error) with which the governing physical laws $P$ are instantiated.

As long as $\epsilon$ remains below the threshold $\epsilon_{\text{DNA}}$ required for the simulated DNA-mediated biochemical and physiological processes to function equivalently to their real-world counterparts, the simulated organism $H'$ will support subjective experience.
Imposing any stricter bound $\epsilon_{\min} < \epsilon_{\text{DNA}}$ specifically for the preservation of $S$ would introduce an additional non-physical constant into the system, violating Axiom 2.

Formally:
\[
\epsilon \leq \epsilon_{\text{DNA}} \implies S(H') \neq \emptyset
\]

In the limit of sufficient fidelity, the simulated structure is axiomatically equivalent to the biological one. Differences in implementation (analog vs. digital, classical vs. quantum, precise vs. approximate) become representational, not ontological.


\subsection*{Falsifiability of the Hypothesis}

The hypothesis is falsifiable in the future when technology advances and DNA simulations can be run with sufficient accuracy for DNA-based organisms.
The effect of pain can be measured just like an effect of physical forces can be measured.
If a DNA simulation $H'$ is constructed and shown to lack $S$, then $A_4$ is invalidated, and the axioms 1--3 collapse.


\subsection*{Conclusions}

The four axioms and the subsequent deductions lead to the following fundamental result:

\begin{principle}{Informational Substrate Principle (ISP)}
{informational-substrate}
The universe is fundamentally a static, timeless, and substrate-independent informational structure.
Subjective experience ($S$), including the perception of time, causality, and physical reality, emerges solely as an internal property of observers encoded within configurations of abstract information. 

No particular physical substrate, dynamical process, or privileged reference frame is required for consciousness. All possible observers correspond to different static arrangements of information, none of which is ontologically primary. The relationship between a ``simulating'' system and a ``simulated'' universe is purely representational, not causal. Time is not an external attribute of reality ($E_{\text{ext}}$) but arises exclusively as an internal experience ($E_{\text{int}}$) within the observer.
\end{principle}

This principle, which we shall refer to as \textbf{ISP} or \textbf{Principle~\ref{principle:informational-substrate}}, serves as the foundational cornerstone for all subsequent derivations in this work.
